\section{Related Work}
\label{sec:related}

\para{Drift detection from neural embeddings.}
Distribution drift detection in deployed ML systems has received growing attention as models are increasingly served in non-stationary environments. \citet{khaki2023uncovering} deploy MMD-based drift detection on BERT embeddings in an Amazon production system, demonstrating 76.9\% correlation between MMD drift scores and model performance degradation over a three-year period. \citet{sodar2023detecting} systematically compare embedding methods (TF-IDF, Doc2Vec, BERT) combined with dimensionality reduction and statistical tests, finding that BERT embeddings with PCA and MMD generally perform best. \citet{greco2024driftlens} propose DriftLens, which uses per-label Fr\'{e}chet distance on embedding distributions for unsupervised drift detection, assuming Gaussian distributions. \citet{kalinke2022mmdew} introduce MMDEW, an online streaming MMD detector with $O(\log^2 t)$ runtime via exponential windows. These methods rely on moment-based or kernel-based statistics that capture changes in location and spread but may miss structural reorganization of the embedding manifold.

\para{TDA for NLP and LLM analysis.}
Topological Data Analysis has found increasing application in natural language processing. \citet{gardinazzi2024persistent} apply zigzag persistence to track topological feature evolution across LLM layers, establishing that topological structure is consistent across architectures and can guide layer pruning. \citet{wei2025shortphd} extend the Persistent Homology Dimension (PHD) metric to detect short LLM-generated texts, exploiting the finding that LLM-generated text exhibits lower PHD than human-written text. A comprehensive survey by \citet{chen2024tdasurvey} documents the growing landscape of TDA applications in NLP, confirming that topological features of text embeddings are discriminative across multiple tasks.

\para{TDA for drift and change detection.}
The most directly related work is \citet{basterrech2024unsupervised}, who combine persistent entropy with topology-preserving dimensionality reduction (Self-Organizing Maps) to detect concept drift. They validate on synthetic MNIST streams grouped by topological homology class (digits with zero, one, or two holes), showing that persistent entropy on SOM-projected data correctly detects injected topological shifts. \citet{bobrowski2022testing} provide theoretical grounding for using Betti numbers as two-sample test statistics, establishing consistency in critical and supercritical regimes of \v{C}ech and Vietoris--Rips complexes.

\para{Positioning of our work.}
Unlike \citet{basterrech2024unsupervised}, who validate only on synthetic image streams, we evaluate on real LLM text embeddings with multiple drift scenarios. Unlike production drift systems~\cite{khaki2023uncovering, greco2024driftlens} that use only moment-based statistics, we include TDA features and identify the specific conditions under which they provide unique detection capability. Our work bridges the TDA-for-NLP literature~\cite{gardinazzi2024persistent, wei2025shortphd} with the drift detection literature~\cite{khaki2023uncovering, sodar2023detecting} by providing the first direct comparison on embedding streams.
